\documentclass[11pt]{article}
\usepackage{amsmath,amssymb,amsthm}
\usepackage{algorithm}
\usepackage[noend]{algpseudocode}
\usepackage{enumitem}
\usepackage{hyperref}
\usepackage{geometry}
\geometry{margin=1in}
\newtheorem{theorem}{Theorem}
\newtheorem{lemma}{Lemma}
\newtheorem{assumption}{Assumption}
\newtheorem{corollary}{Corollary}
\newcommand{\argmin}{\operatorname*{arg\,min}}
\newcommand{\D}{\mathcal{D}}
\newcommand{\R}{\mathbb{R}}
\newcommand{\E}{\mathbb{E}}

\begin{document}

%=====================================================================
\section*{Algorithm Name}
\textbf{Mirror Descent with Bregman Proximal Step}
%=====================================================================

%=====================================================================
\section*{Mathematical Formulation}
Let $f:\R^{n}\rightarrow\R$ be a convex differentiable objective and let 
$h:\R^{n}\rightarrow\R$ be a \emph{prox‐function} that is $\sigma$–strongly convex
with respect to a norm $\|\cdot\|$:
\[
h(y)\ge h(x)+\langle\nabla h(x),y-x\rangle+\frac{\sigma}{2}\|y-x\|^{2},
\qquad\forall x,y\in\R^{n}.
\]
The associated Bregman divergence is
\[
D_{h}(y,x):=h(y)-h(x)-\langle\nabla h(x),y-x\rangle\ge
\frac{\sigma}{2}\|y-x\|^{2}. \tag{1}
\]

Given a stepsize sequence $\{\eta_{k}\}_{k\ge0}$ with $\eta_{k}>0$, the
\emph{mirror descent} iterate is defined as
\[
x_{k+1}
= \arg\min_{x\in\R^{n}}\Big\{
\langle \nabla f(x_{k}),x\rangle
+\frac{1}{\eta_{k}} D_{h}(x,x_{k})
\Big\}. \tag{2}
\]

The optimality condition of (2) reads
\[
\nabla h(x_{k+1}) = \nabla h(x_{k}) - \eta_{k}\,\nabla f(x_{k}). \tag{3}
\]

%=====================================================================
\section*{Pseudocode}
\begin{algorithm}[H]
\caption{Mirror Descent}
\begin{algorithmic}[1]
\State \textbf{Input:} initial point $x_{0}\in\R^{n}$, stepsize rule $\{\eta_{k}\}$,
prox‑function $h$ (strongly convex).
\For{$k=0,1,2,\dots,T-1$}
    \State Compute gradient $g_{k}\gets \nabla f(x_{k})$.
    \State Compute 
    $x_{k+1}\gets\arg\min_{x}\big\{\langle g_{k},x\rangle+\frac{1}{\eta_{k}}D_{h}(x,x_{k})\big\}$.
    \Comment{equivalently solve (3): $\nabla h(x_{k+1})=\nabla h(x_{k})-\eta_{k}g_{k}$}
\EndFor
\State \textbf{Output:} $x_{T}$ (or the averaged point $\bar x_{T}=\frac{1}{T}\sum_{k=0}^{T-1}x_{k}$).
\end{algorithmic}
\end{algorithm}

%=====================================================================
\section*{Dry Run Example}
Consider the scalar problem
\[
f(w)=(w-3)^{2},\qquad h(w)=\tfrac12 w^{2}\;( \text{so } D_{h}(u,v)=\tfrac12 (u-v)^{2}),
\]
with stepsize $\eta_{k}= \eta =0.1$ and $w_{0}=0$.

\begin{enumerate}[label=Iteration \arabic*:, leftmargin=*, align=left]
\item \textbf{Gradient:} $\nabla f(w_{0}) = 2(w_{0}-3) = -6$.
      \\
      \textbf{Update:} 
      \[
      w_{1}= \arg\min_{w}\Big\{-6w+\frac{1}{0.1}\tfrac12 (w-0)^{2}\Big\}
      = w_{0}-\eta\,\nabla f(w_{0}) = 0-0.1(-6)=0.6 .
      \]
      \textbf{Objective:} $f(w_{1})=(0.6-3)^{2}=5.76$.

\item \textbf{Gradient:} $\nabla f(w_{1}) = 2(0.6-3) = -4.8$.
      \\
      \textbf{Update:}
      \[
      w_{2}= w_{1}-\eta\,\nabla f(w_{1})=0.6-0.1(-4.8)=1.08 .
      \]
      \textbf{Objective:} $f(w_{2})=(1.08-3)^{2}=3.6864$.

\item \textbf{Gradient:} $\nabla f(w_{2}) = 2(1.08-3) = -3.84$.
      \\
      \textbf{Update:}
      \[
      w_{3}= w_{2}-\eta\,\nabla f(w_{2})=1.08-0.1(-3.84)=1.464 .
      \]
      \textbf{Objective:} $f(w_{3})=(1.464-3)^{2}=2.366 .
      \]
\end{enumerate}
The sequence follows the classical gradient step because the chosen
prox‑function is Euclidean; for a non‑Euclidean $h$ the update would differ.

%=====================================================================
\section*{Assumptions}
\begin{assumption}[Convexity and Smoothness Relative to $h$]
\label{ass:convex-smooth}
The objective $f$ satisfies:
\begin{enumerate}[label=(\alph*)]
\item $f$ is convex:
\[
f(y)\ge f(x)+\langle\nabla f(x),y-x\rangle,\qquad\forall x,y.
\]
\item $f$ is $L$–smooth \emph{relative} to $h$:
\[
f(y)\le f(x)+\langle\nabla f(x),y-x\rangle+L\,D_{h}(y,x),
\qquad\forall x,y. \tag{4}
\]
\end{enumerate}
\end{assumption}

\begin{assumption}[Strong Convexity of the Prox‑function]
\label{ass:h-strong}
The prox‑function $h$ is $\sigma$–strongly convex w.r.t. $\|\cdot\|$,
i.e., inequality (1) holds with constant $\sigma>0$.
\end{assumption}

\begin{assumption}[Stepsize]
\label{ass:stepsize}
The stepsizes are constant,
\[
\eta_{k}\equiv\eta =\frac{1}{L},
\]
or more generally $0<\eta_{k}\le\frac{1}{L}$ for all $k$.
\end{assumption}

%=====================================================================
\section*{Main Convergence Theorem}
\begin{theorem}[Convergence of Mirror Descent]
\label{thm:md}
Assume \ref{ass:convex-smooth}–\ref{ass:stepsize}. Let $x^{\star}\in\arg\min f$
and define the Bregman distance $D_{h}(x^{\star},x_{0})$.
Then for any $T\ge1$ the iterates generated by \eqref{2} satisfy
\[
f(\bar x_{T})-f(x^{\star})
\;\le\;
\frac{L}{\sigma}\,\frac{D_{h}(x^{\star},x_{0})}{T},
\qquad
\bar x_{T}:=\frac{1}{T}\sum_{k=0}^{T-1}x_{k}.
\tag{5}
\]
Consequently, the optimality gap decays as $O(1/T)$.
\end{theorem}

\begin{theorem}[Linear Rate under Strong Convexity of $f$]
\label{thm:strong}
If, in addition, $f$ is $\mu$–strongly convex (ordinary, not relative),
\[
f(y)\ge f(x)+\langle\nabla f(x),y-x\rangle+\frac{\mu}{2}\|y-x\|^{2},
\]
then with a constant stepsize $\eta=\frac{1}{L+\mu}$,
\[
f(x_{k})-f(x^{\star})\le\Big(1-\frac{\mu}{L+\mu}\Big)^{k}
\big(f(x_{0})-f(x^{\star})\big).
\tag{6}
\]
\end{theorem}

%=====================================================================
\section*{Theoretical Properties}
\begin{itemize}
\item \textbf{Stability:} The update \eqref{3} is a contraction in the dual
space $\nabla h(\cdot)$ because of strong convexity of $h$.
\item \textbf{Hyperparameter Sensitivity:}
Stepsize must satisfy $\eta\le1/L$; larger $\eta$ may violate the
smoothness inequality (4) and break the $O(1/T)$ guarantee.
\item \textbf{Comparison to Gradient Descent:}
When $h(w)=\tfrac12\|w\|^{2}$, $D_{h}$ reduces to Euclidean distance and
\eqref{2} becomes ordinary gradient descent. For non‑Euclidean $h$,
mirror descent can exploit geometry (e.g., KL divergence for simplex
constraints) and often yields tighter constants $\frac{L}{\sigma}$.
\end{itemize}

%=====================================================================
\section*{Discussion}
The theorem shows that only the relative smoothness constant $L$ and the
strong convexity constant $\sigma$ of the prox‑function appear in the
rate.  By choosing $h$ that matches the feasible set (e.g., negative
entropy on the probability simplex) one can obtain $\sigma$ close to $1$
and dramatically improve practical convergence compared with plain
gradient descent.

%=====================================================================
\section*{Preliminaries (Definitions and Lemmas)}
\begin{lemma}[Three‑Point Identity]
\label{lem:three-point}
For any $x,y,z\in\R^{n}$,
\[
\langle \nabla h(z)-\nabla h(y),x-y\rangle
= D_{h}(x,y)+D_{h}(y,z)-D_{h}(x,z). \tag{7}
\]
\end{lemma}
\begin{proof}
Direct expansion of the definition of $D_{h}$ yields (7). \hfill\qedhere
\end{proof}

\begin{lemma}[Descent Lemma (Relative Smoothness)]
\label{lem:rel-smooth}
If $f$ satisfies (4) then for any $x,y$,
\[
f(y)\le f(x)+\langle\nabla f(x),y-x\rangle+L D_{h}(y,x). \tag{8}
\]
\end{lemma}
\begin{proof}
Lemma statement is exactly the definition of relative smoothness. \hfill\qedhere
\end{proof}

%=====================================================================
\section*{Main Proof (Step‑by‑Step Derivation)}
We prove Theorem~\ref{thm:md}.  Throughout the proof let $x^{\star}\in
\arg\min f$ be fixed.

\begin{enumerate}[label=[Step \arabic*], leftmargin=*, align=left]
\item \textbf{Optimality condition of the proximal subproblem.}\\
From the definition \eqref{2},
\[
0\in \nabla f(x_{k})+\frac{1}{\eta_{k}}\big(\nabla h(x_{k+1})-\nabla h(x_{k})\big).
\]
Rearranging gives exactly \eqref{3}:
\[
\nabla h(x_{k+1})=\nabla h(x_{k})-\eta_{k}\nabla f(x_{k}). \tag{9}
\]

\item \textbf{Apply the three‑point identity with $x=x^{\star}$,
$y=x_{k}$, $z=x_{k+1}$.}\\
Using Lemma~\ref{lem:three-point} we obtain
\[
\langle \nabla h(x_{k+1})-\nabla h(x_{k}),x^{\star}-x_{k}\rangle
= D_{h}(x^{\star},x_{k})+D_{h}(x_{k},x_{k+1})-D_{h}(x^{\star},x_{k+1}). \tag{10}
\]

\item \textbf{Replace the left‑hand side using (9).}\\
From (9),
\[
\langle -\eta_{k}\nabla f(x_{k}),x^{\star}-x_{k}\rangle
= D_{h}(x^{\star},x_{k})+D_{h}(x_{k},x_{k+1})-D_{h}(x^{\star},x_{k+1}). \tag{11}
\]

\item \textbf{Rearrange (11) to isolate the inner product.}\\
Dividing by $\eta_{k}>0$,
\[
\langle \nabla f(x_{k}),x_{k}-x^{\star}\rangle
=\frac{1}{\eta_{k}}\big[ D_{h}(x^{\star},x_{k})-D_{h}(x^{\star},x_{k+1})
- D_{h}(x_{k+1},x_{k})\big]. \tag{12}
\]

\item \textbf{Use convexity of $f$.}\\
Convexity yields
\[
f(x_{k})-f(x^{\star})\le \langle \nabla f(x_{k}),x_{k}-x^{\star}\rangle.
\tag{13}
\]

\item \textbf{Combine (12) and (13).}\\
\[
f(x_{k})-f(x^{\star})
\le \frac{1}{\eta_{k}}\big[ D_{h}(x^{\star},x_{k})-D_{h}(x^{\star},x_{k+1})
- D_{h}(x_{k+1},x_{k})\big]. \tag{14}
\]

\item \textbf{Drop the non‑negative term $D_{h}(x_{k+1},x_{k})$.}\\
Since $D_{h}\ge0$,
\[
f(x_{k})-f(x^{\star})
\le \frac{1}{\eta_{k}}\big[ D_{h}(x^{\star},x_{k})-D_{h}(x^{\star},x_{k+1})\big].
\tag{15}
\]

\item \textbf{Sum (15) from $k=0$ to $T-1$.}\\
Telescoping gives
\[
\sum_{k=0}^{T-1}\big(f(x_{k})-f(x^{\star})\big)
\le \frac{1}{\eta}\,D_{h}(x^{\star},x_{0}), \tag{16}
\]
where we used the constant stepsize $\eta_{k}\equiv\eta$ (Assumption~\ref{ass:stepsize}).

\item \textbf{Relate the average iterate to the sum.}\\
Define $\bar x_{T}:=\frac{1}{T}\sum_{k=0}^{T-1}x_{k}$. By convexity of $f$,
\[
f(\bar x_{T})-f(x^{\star})\le \frac{1}{T}\sum_{k=0}^{T-1}\big(f(x_{k})-f(x^{\star})\big).
\tag{17}
\]

\item \textbf{Combine (16) and (17).}\\
\[
f(\bar x_{T})-f(x^{\star})
\le \frac{1}{\eta T}\,D_{h}(x^{\star},x_{0}). \tag{18}
\]

\item \textbf{Insert the admissible stepsize $\eta=1/L$ and the strong convexity
bound (1).}\\
Since $D_{h}(x^{\star},x_{0})\le \frac{1}{2\sigma}\|x^{\star}-x_{0}\|^{2}$,
\[
f(\bar x_{T})-f(x^{\star})
\le \frac{L}{\sigma}\,\frac{D_{h}(x^{\star},x_{0})}{T},
\tag{19}
\]
which is exactly the claim \eqref{5}. \hfill\qedhere
\end{enumerate}

%=====================================================================
\section*{Rate Derivation (With Constants)}
From \eqref{19} we read the explicit constant:
\[
C:=\frac{L}{\sigma}D_{h}(x^{\star},x_{0}),
\qquad\text{so}\qquad
f(\bar x_{T})-f(x^{\star})\le \frac{C}{T}.
\]
If $h$ is Euclidean ($h(x)=\tfrac12\|x\|^{2}$) then $\sigma=1$ and
$D_{h}(x^{\star},x_{0})=\tfrac12\|x^{\star}-x_{0}\|^{2}$, recovering the
standard $O(L\|x_{0}-x^{\star}\|^{2}/T)$ rate of gradient descent.

%=====================================================================
\section*{Special Cases}
\begin{itemize}
\item \textbf{Euclidean Setting.} $h(x)=\frac12\|x\|^{2}$ yields
$D_{h}(y,x)=\frac12\|y-x\|^{2}$ and $\sigma=1$. The algorithm reduces to
gradient descent with step $1/L$ and Theorem~\ref{thm:md} coincides with
the classical $O(1/T)$ bound.
\item \textbf{Probability Simplex.} Choose $h(x)=\sum_{i}x_{i}\log x_{i}$,
the negative entropy. Then $D_{h}$ is the Kullback–Leibler divergence,
$\sigma=1$, and mirror descent becomes the exponentiated gradient method.
The same $O(1/T)$ guarantee holds but now the iterates stay in the simplex
without projection.
\item \textbf{Strongly Convex $f$.} Under the additional ordinary strong
convexity of $f$, Theorem~\ref{thm:strong} gives a linear convergence rate,
again with the same constant stepsize $\eta=1/(L+\mu)$.
\end{itemize}

%=====================================================================
\section*{References}
\begin{enumerate}
\item A. Beck and M. Teboulle, “Mirror descent and nonlinear projected
gradient methods,” \emph{Operations Research Letters}, 2003.
\item Y. Nesterov, “Primal–dual subgradient methods for convex problems,”
\emph{Mathematical Programming}, 2009.
\end{enumerate}

\end{document}